% Root mathematical contribution. Finding IDs identify deductions checked in this run.
\section{The exact abelian Artin--Tate consequence}
\label{sec:root-artin-tate}

We use the cohomological convention $\mathbb Q(1)$ of weight $-2$,
so that a geometric Frobenius at a good prime acts on its $\ell$-adic
realization by $p^{-1}$.  Coefficients may be extended to a number field
containing the values of the finite characters involved.  A rational
statement is recovered by grouping coefficient-conjugate characters.

\paragraph{Finding BR-AT-01 (the full shifted-zero consequence).}
Let $M$ be a finite direct sum of rank-one pure motives
\[
 M=\bigoplus_{\chi,n} (\chi\otimes\mathbb Q(n))^{\oplus m_{\chi,n}},
 \qquad m_{\chi,n}\in\mathbb Z_{\geq0},\quad n\in\mathbb Z,
\]
where $\chi$ is a finite-order character of
$\operatorname{Gal}(\overline{\mathbb Q}/\mathbb Q)$ with abelian image.
Use the primitive Dirichlet character associated with each $\chi$ for
its full Artin $L$-function.  Then
\[
 L(M,s)=\prod_{\chi,n}L(\chi,s+n)^{m_{\chi,n}}.
\]
Under the adopted Alp\"oge theorem, the nontrivial zeros contributed by
the summand of twist $n$ have real part $\tfrac12-n$.
These zero multiplicities add in the displayed product: none is cancelled
by a pole of another factor.  For a pure summand of weight $w=-2n$, this
is the expected line $\Re s=(w+1)/2$.

For a mixed object $V$ whose weight-graded pieces have this form, and a
finite set $S$ containing every ramified prime, the precise assertion is
\[
 L^S(V,s)=L^S(\operatorname{gr}^W V,s)
 =\prod_{\chi,n}L^S(\chi,s+n)^{m_{\chi,n}}.
\]
If Euler factors at primes in $S$ are deleted, the resulting extra zeros
lie on the lines $\Re s=-n$; they must not be called violations of the
nontrivial-zero assertion.  At ramified primes, no general multiplicativity
of inertia invariants for mixed extensions is being assumed.

\begin{proof}
At an unramified prime $p$, the eigenvalue for $\chi\otimes\mathbb Q(n)$
is $\chi(p)p^{-n}$, giving the Euler factor
$(1-\chi(p)p^{-s-n})^{-1}$.  The identity of the convergent Euler
products follows in a right half-plane and extends to their meromorphic
continuations.  This also proves the direct-sum formula.  On replacing
$z=s+n$, the strip $0<\Re z<1$ becomes
$-n<\Re s<1-n$; the adopted Dirichlet theorem places every zero there
on $\Re s=1/2-n$.  A primitive nonprincipal Dirichlet $L$-function is
entire; a principal primitive factor is $\zeta(s+n)$ and has its only
pole at $s=1-n$.  Such poles have integral real part, while the indicated
nontrivial zeros have half-integral real part.  Hence they cannot cancel.

For a filtration preserved by Frobenius, choose a basis respecting the
filtration.  Its matrix is block upper triangular; the determinant of
$1-p^{-s}\operatorname{Frob}_p$ is the product of the determinants of
the diagonal blocks.  At every $p\notin S$ this gives the mixed formula.
Deleting a nonzero local rank-one factor multiplies the full $L$-function
by $1-\alpha p^{-s-n}$ with $|\alpha|=1$.  Its zeros satisfy
$p^{-\Re s-n}=1$, hence $\Re s=-n$.  A ramified one-dimensional
character with zero inertia-invariant space contributes the factor $1$,
so there is nothing to delete in that case.  This proves all the assertions.
\end{proof}

The elementary Euler-factor argument is included because it specifies the
mixed, twisted, and incomplete forms not stated by the Lean declarations.
It does not assert that every general mixed-motivic $L$-function is
determined at every bad prime by its semisimplification.  In particular,
calling an object ``mixed Artin--Tate'' does not make its extension class
part of a Dirichlet character.

\paragraph{Archimedean convention.}
For $\chi(-1)=(-1)^a$, $a\in\{0,1\}$, the factor of
$\chi\otimes\mathbb Q(n)$ is $\Gamma_{\mathbb R}(s+n+a)$, where
$\Gamma_{\mathbb R}(s)=\pi^{-s/2}\Gamma(s/2)$.  Its real Frobenius
acts as $\chi(-1)(-1)^n$.  Both the Hodge type $(-n,-n)$ and that
sign are needed to recover the archimedean factor.  The sign alone does
not recover the weight: $\mathbb Q(0)$ and $\mathbb Q(2)$ both have
trivial real Frobenius.  The corrected real-Frobenius classification at
cote 18/038 is treated with the original text below.

\subsection{Mixed coefficients: a rigorous obstruction to reconstructing extensions}

Deligne--Goncharov (2005), \S1.1, equations (1.1.2), (1.1.4), and the
proof of Proposition 1.7, construct the Kummer extension $K(x)$ and
identify
\[
 \operatorname{Ext}^1_{\mathrm{MT}(\mathbb Q)}
   (\mathbb Q(0),\mathbb Q(1))=\mathbb Q^\times\otimes_{\mathbb Z}\mathbb Q.
\]
These exact locators are in the supplied file \texttt{Sec1ccc.tex}; this
identification is a cited theorem, not a new result here.

\paragraph{Finding BR-MIX-01 (the Kummer extension survives identical
good Euler factors).}
There are nonisomorphic mixed Tate motives with the same weight-graded
object and the same Euler polynomials at every prime of simultaneous
good reduction.  Explicitly,
\[
 0\longrightarrow\mathbb Q(1)\longrightarrow K(2)
 \longrightarrow\mathbb Q(0)\longrightarrow0
\]
is nonsplit, whereas $\mathbb Q(0)\oplus\mathbb Q(1)$ is split.  For
$S$ containing $2$ and the realization prime $\ell$, both have
\[
 L^S(s)=\zeta^S(s)\zeta^S(s+1).
\]
Thus even complete knowledge of these good Euler factors, their
meromorphic continuation, and all their zeros does not reconstruct
the mixed extension.  This refutes that particular reconstruction
rule; it is not a logical independence theorem about RH.

\begin{proof}
The $2$-adic valuation gives a homomorphism
$v_2:\mathbb Q^\times\to\mathbb Z$ and hence a linear map
$\mathbb Q^\times\otimes\mathbb Q\to\mathbb Q$ taking $2\otimes1$
to $1$.  Consequently the class of $2$ is nonzero, and the cited
Ext-identification makes its extension nonsplit.  An isomorphism to
the split object would preserve the unique weight filtration and
give a splitting, a contradiction.  By Deligne--Goncharov,
Proposition 1.7 and its proof, this Kummer object is unramified away
from $2$ (and the $\ell$-adic realization is used away from $\ell$).
For $p\notin S$ its Frobenius matrix in a weight-adapted basis is
upper triangular with diagonal entries $p^{-1},1$.  Its Euler
polynomial is therefore $(1-T)(1-p^{-1}T)$, the polynomial of the
split object.  Taking their Euler products proves the identity.
\end{proof}

The same proof works for every $x\in\mathbb Q^\times$ with a nonzero
valuation somewhere, after deleting the primes dividing $x$.
The unipotent radical of the mixed Tate motivic Galois group, its
extension classes, and its period matrix are additional information.
This is a concrete reason why a theorem about the zero divisor of
these $L$-functions does not by itself compute the groups in step 17
of \emph{Remark}.  In particular the tentative footnote there gives
neither the degrees nor the primary parts required for a $K$-group
vanishing statement.  The assertion that \emph{all} positive-degree
$K$-groups of $\mathbb Z$ vanish is false already for
$K_1(\mathbb Z)=\{\pm1\}$; it is not an assertion made by that footnote.

\section{Cote 1: what the Hilbert-space trace theorem settles}
\label{sec:root-trace}

The supplied scans 1/079--080 contain the Fredholm determinant estimates
and the trace/eigenvalue formula; 1/111--112 contain the genus-zero
argument.  They do not mention $\xi$, $\Xi$, a Riemann-zero operator,
or RH.  Those are applications in the supplied SA23 intake, \S4.1,
and in this work.  An exact attribution therefore distinguishes the
original analytic theorem from its later arithmetic use.

Write
\[
 \xi(s)=\tfrac12s(s-1)\pi^{-s/2}\Gamma(s/2)\zeta(s),\qquad
 \Xi(t)=\xi(\tfrac12+it),\qquad F(w)=\frac{\Xi(\sqrt w)}{\Xi(0)}.
\]
Evenness of $\Xi$ makes $F$ single-valued and entire.  The classical
Hadamard product and zero-counting facts used here are Titchmarsh--
Heath-Brown (1986), \S2.12 (the $\xi$ product), Theorems 9.3--9.4,
and equation (9.4.3), p.~214 (the Riemann--von Mangoldt count).
We use their statements, not a
new proof of those analytic inputs.

\paragraph{Finding BR-TR-01 (a positive trace-class realization, and its limit).}
There is a diagonal trace-class operator $B$ with
\[
 F(w)=\det(I-wB).
\]
It can be constructed before RH by taking one zero $t_j$ from each
pair $\{t,-t\}$ of $\Xi$, counted with multiplicity, and putting
$B=\operatorname{diag}(t_j^{-2})$ on $\ell^2$.
For this normal realization, RH is equivalent to $B\geq0$.
Consequently the adopted RH gives such a positive realization.
Every trace-class operator $T$ with $\det(I-wT)=F(w)$ has the same
nonzero eigenvalues with algebraic multiplicity and satisfies
\[
 \operatorname{Tr}(T^k)=\sum_jt_j^{-2k}>0\quad(k\geq1)
\]
under RH.  It is false that every such $T$ is positive or self-adjoint.
Moreover, no trace-class operator has $\Xi(t)$ itself as an ordinary
Fredholm determinant in an affine variable $t$.

\begin{proof}
The nontrivial zeta zeros have counting function $O(T\log(T+2))$.
Partial summation, or a dyadic decomposition, gives
$\sum_j|t_j|^{-2}<\infty$; finitely many bounded zeros cause no issue
because $\Xi(0)\ne0$.  Thus the displayed diagonal operator is trace
class.  The entire even function $\Xi$ has order one, so $F$ has
order at most $1/2$.  Hadamard factorization in order less than one,
together with $F(0)=1$, gives the genus-zero product
$F(w)=\prod_j(1-wt_j^{-2})$, which is precisely its Fredholm determinant.

If RH holds, all $t_j$ are real and nonzero; thus every diagonal
entry of $B$ is positive.  Conversely $B\geq0$ implies
$t_j^{-2}>0$, so $t_j$ is real.  Indeed, for a nonzero complex number
$t$, $t^2$ positive real forces $t$ real; the purely imaginary
alternative has negative square.  Hence every zero of $\xi(s)$
has real part $1/2$.

For a general trace-class $T$, the Fredholm zero theorem identifies
its nonzero eigenvalues and their algebraic multiplicities with the
reciprocals of the zeros of its determinant.  The trace theorem then
gives the power-trace formula; the sum converges absolutely.  To see
why positivity of $T$ does not follow, start with $B$ under RH and
replace its first two diagonal entries $b_1,b_2$ by the block
$\left(\begin{smallmatrix}b_1&1\\0&b_2\end{smallmatrix}\right)$.
This is a finite-rank perturbation, remains trace class, has the
same determinant, and is not self-adjoint in the given Hilbert
inner product.  It therefore is not a positive operator there.

Finally suppose $\Xi(t)=c\det(I+(t-a)T)$ for trace-class $T$ and
$c\ne0$.  Its zeros would be $a-\lambda_j^{-1}$, with the nonzero
eigenvalues $\lambda_j$ of $T$ counted algebraically.  The trace-class
eigenvalue inequality forces $\sum_j|\lambda_j|<\infty$.
But the Riemann--von Mangoldt asymptotic gives divergent sum of the
reciprocals of the distances of the $\Xi$ zeros from any fixed $a$.
For example, in dyadic height intervals the count is of order
$2^k k$, whereas every such distance is at most $C2^k$ for large
$k$ (the imaginary parts of the $t$-zeros are bounded by $1/2$).
The sum over that interval is at least a positive constant times
$k$.  This contradicts the required summability.  A zero at $a$
would already contradict the nonzero determinant value at $t=a$.
\end{proof}

The construction uses the zero set; it does not solve the stronger
geometric Hilbert--P\'olya problem of producing the operator from
primes, correspondences, or a cohomology theory independently of its
prescribed spectrum.  The result settles the exact positive normal
realization and its trace inequalities, not that additional
construction problem.

\paragraph{Finding BR-TR-02 (a strictly positive moment hierarchy).}
Under RH put $b_j=t_j^{-2}>0$ and
$m_r=\sum_j b_j^{r+1}$ for $r\geq0$.  For every integer $d\geq0$,
both Hankel matrices $(m_{a+b})_{0\leq a,b\leq d}$ and
$(m_{a+b+1})_{0\leq a,b\leq d}$ are positive definite.
For $x\geq0$,
\[
 H(x)=-\frac{F'(-x)}{F(-x)}
      =\operatorname{Tr}\bigl(B(I+xB)^{-1}\bigr)
      =\sum_j\frac{b_j}{1+xb_j},
 \qquad (-1)^r H^{(r)}(x)>0\quad(r\geq0).
\]
These are quantitative inequalities in the cote 1 trace framework,
not assertions that its unspecified geometric operators are positive.

\begin{proof}
The positive measure $\mu=\sum_jb_j\delta_{b_j}$ is finite because
$B$ is trace class, and it has infinitely many distinct support points:
$\Xi$ has infinitely many zeros and finite multiplicity at each point.
For a nonzero polynomial $P(u)=\sum_{a=0}^d c_au^a$,
\[
 \sum_{a,b=0}^d\overline{c_a}c_bm_{a+b}
   =\int |P(u)|^2\,d\mu(u)>0.
\]
A nonzero polynomial cannot vanish on infinitely many distinct points.
Inserting an additional factor $u>0$ proves the second positive
definiteness assertion.  The logarithmic derivative follows from
the normally convergent product for $F$; for $x\geq0$ all its
denominators are positive.  On compact subintervals of $[0,\infty)$,
termwise differentiation is dominated by
$r!\|B\|^r\sum_jb_j$.  Thus
\[
 (-1)^rH^{(r)}(x)
 =r!\sum_j\frac{b_j^{r+1}}{(1+xb_j)^{r+1}}>0.
\]
This proves convergence, the identities, and strictness.
\end{proof}

\section{Weights, monodromy, and the additional geometric input}
\label{sec:root-weights}

Deligne (1980), \emph{La conjecture de Weil II}, Theorem 1.8.4,
states that for a punctually pure lisse sheaf of weight $w$ on the
punctured smooth curve over a finite field, the $j$th graded piece
of its local monodromy filtration is pure of weight $w+j$.
This is a theorem with a geometric-purity hypothesis, not an
implication from a list of zeros of $\zeta(s)$.
The mixed-characteristic question for $H^i$ of a smooth proper
variety over a local field is the corresponding assertion that
$\operatorname{gr}_j^M H^i$ has weight $i+j$.
An exact modern statement attributing it to Deligne is Scholze (2012),
\emph{Perfectoid Spaces}, Conjecture 9.3; Theorem 9.6 proves the
set-theoretic complete-intersection case in smooth projective toric
varieties.  No general mixed-characteristic proof is contained in
the supplied RH declarations.

\paragraph{Finding BR-WM-01 (a determinant and the Tate covariance do not
determine monodromy purity).}
For every prime power $q>1$ there are two Weil--Deligne linear
algebra data with the same semisimple Frobenius and the same
Frobenius determinant, both satisfying $FNF^{-1}=q^{-1}N$, for
which the monodromy-purity statement centered at weight $1$ has
different truth values.

\begin{proof}
Take a two-dimensional vector space with basis $e_0,e_2$ and
$Fe_0=e_0$, $Fe_2=qe_2$.  In the first datum put
$Ne_2=e_0$, $Ne_0=0$; in the second put $N=0$.
The covariance is checked on the basis and holds in both cases.
Their common Frobenius determinant is $(1-T)(1-qT)$.
For the first datum the monodromy filtration centered at zero is
$M_{-2}=0$, $M_{-1}=M_0=\langle e_0\rangle$, $M_1=V$.
Its graded pieces in degrees $-1$ and $1$ have eigenvalues $1$
and $q$, hence weights $0$ and $2$, which are $1-1$ and $1+1$.
For the second datum the monodromy filtration is concentrated in
degree zero, so its only graded piece is $V$, with weights $0,2$;
it is not pure of weight $1$.  This proves the claim.
\end{proof}

These data refute the proposed \emph{linear-algebra deduction} from
Frobenius roots and covariance.  They are not a counterexample to
Deligne's geometric conjecture: no realization as the cohomology
of a smooth proper variety of the required weight is asserted for
the second datum.  Also, the local Euler factor defined on
$\ker N$ differs between the two examples.  It would be incorrect
to call these examples indistinguishable by every ramified local
$L$-factor.  What fails is precisely the attempted passage from
a determinant on $V$ to its monodromy-graded purity.

\subsection{The tame triangle and the nodal universal curve}

The intake SA20, row C.008, attributes its triangle to cote 122;
the original scan for that row is not among the supplied files.
Its finite $\mathbb Z/\ell^a$ coefficients and twists type the
duality maps.  Numerical Frobenius weights require a rational
$\ell$-adic realization and a purity theorem.  After that passage,
$H^1(I_t,\mathbb Q_\ell)=\mathbb Q_\ell(-1)$ has weight $2$;
the monodromy logarithm lowers weights by $2$.  These facts explain
the appearance of adjacent twists, but do not prove the
mixed-characteristic monodromy-purity assertion.

There is a further distinction in the RESIDUE intake, H.3--H.6:
the ordinary cohomology of a singular nodal fibre is not the same
filtered object as the nearby smooth cohomology with its full
monodromy filtration.  The following elementary calculation makes
the distinction exact.

\paragraph{Finding BR-NODE-01 (the weight-zero part belongs to the dual graph).}
Let $C/\mathbb F_q$ be a geometrically connected proper nodal curve,
with $\ell\ne\operatorname{char}\mathbb F_q$, whose geometric
components, nodes, and two branches at every node are all defined over
$\mathbb F_q$.  Write $\widetilde C=\coprod_v C_v$ for the
normalization and $\Gamma$ for the dual graph.  There is a
Frobenius-equivariant exact sequence
\[
 0\longrightarrow H^1(\Gamma,\mathbb Q_\ell)
 \longrightarrow H^1(C_{\overline{\mathbb F}_q},\mathbb Q_\ell)
 \longrightarrow\bigoplus_vH^1(C_v,\mathbb Q_\ell)
 \longrightarrow0.
\]
The left term has weight zero and the right term weight one.
In particular a split rational curve with one ordinary node has
$H^1\simeq\mathbb Q_\ell$ of weight zero.  Its smooth genus-one
nearby fibre, if chosen in a semistable degeneration, also has the
weight-two monodromy quotient; that quotient is not an additional
piece of the singular fibre's ordinary $H^1$.

\begin{proof}
For the finite normalization map $\nu$, the exact sequence of
\'etale sheaves is
\[
 0\longrightarrow\mathbb Q_\ell
 \longrightarrow\nu_*\mathbb Q_\ell
 \longrightarrow\bigoplus_{e\in E(\Gamma)}\mathbb Q_{\ell,e}
 \longrightarrow0.
\]
Away from the nodes the first map is an isomorphism.  At a node
its stalk is the diagonal inclusion $\mathbb Q_\ell\to
\mathbb Q_\ell^2$, whose cokernel is one-dimensional; choosing
an orientation identifies the quotient with the difference of
branch values.  Finite pushforward and cohomology give the usual
normalization cohomology.  The first part of the long exact
sequence is consequently
\[
 0\to\mathbb Q_\ell\to\mathbb Q_\ell^{V(\Gamma)}
 \xrightarrow{\partial}\mathbb Q_\ell^{E(\Gamma)}
 \to H^1(C)\to\bigoplus_vH^1(C_v)\to0.
\]
The cokernel of the incidence map $\partial$ is graph cohomology.
All vertices and branches are rational, so Frobenius acts trivially
on this graph complex, giving weight zero.  Each $C_v$ is a smooth
proper curve; its $H^1$ has weight one by the finite-field Weil
theorem.  If $C$ is rational with one node, its graph has one
loop and its normalization has no $H^1$, proving the special case.
For a semistable genus-one smoothing, use the cited Tate-uniformization
calculation (Ochiai, 2010, \S4, p.~4, displayed Tate-module and
$H^1$ exact sequences; Berger, 2004, \S II.4.1--II.4.3 gives the
period calculation): the nonzero rank-one monodromy pairs the
weight-zero subspace with a weight-two quotient, exactly as in the
first datum of BR-WM-01.
The nearby fibre has dimension two in $H^1$, whereas the singular
fibre has dimension one.  Thus the extra quotient cannot be
ordinary $H^1$ of the singular curve.
\end{proof}

The already established decomposition theorem concerns the pure
\emph{complex} pushed forward by a proper morphism.  It does not
assert that every ordinary stalk cohomology group of that complex
is pure of its cohomological degree at a singular stratum.  This
distinction prevents a false purity claim for the boundary curves
of $\overline{\mathcal M}_{g,n}$.  The supplied entrance-category
claims H.3--H.6 remain background at the level at which their intake
checked them; their missing original proofs have not been silently
promoted to checks made in this run.

\subsection{The standard conjectures: which additional gate remains}

Grothendieck's \emph{Standard conjectures on algebraic cycles}
(1969), \S2, p.~195, asks for algebraicity of the K\"unneth
projectors (C); \S3, pp.~195--197, asks for the Lefschetz
standard conjectures (A and B); \S4, pp.~197--198, asks for
positivity of the primitive algebraic intersection form (Hdg).
The implication to the finite-field Weil conjectures is stated in
\S4, Remark (2), p.~197; the relation to numerical equivalence
is in Remarks (3)--(4), pp.~197--198.
There is an essential positive consequence of finite-field purity:
Katz--Messing (1974), Theorem 2(1), pp.~76--77, constructs algebraic
K\"unneth projectors over finite fields.  The distinct weights make the
Frobenius characteristic polynomials in different degrees coprime;
Chinese-remainder polynomials in the algebraic Frobenius correspondence
then act as the cohomological projectors.  Theorem 1 of that paper
transfers the required Frobenius characteristic polynomials to crystalline
cohomology under its stated axioms.  These are already established
consequences of the finite-field theorem, and a complete new proof of
that theorem would give another route to them.  They do not, by
themselves, construct the inverse Lefschetz correspondence or prove
positivity on algebraic primitive cycles in general.  The classical
number-field Dirichlet assertion does not add those correspondences.

For the abelian Artin--Tate subcategory the relevant structures are
already explicit: Artin projectors are obtained from finite
correspondences and character idempotents, and Tate objects have
their standard duality.  BR-AT-01 settles the global zero-location
assertion for their $L$-functions.  For arbitrary smooth projective
varieties the current gate is still an algebraic correspondence
realizing the inverse Lefschetz operator, together with the
required positive pairing on algebraic primitive classes.  This
work supplies neither that correspondence nor a counterexample to
the standard conjectures.  The unresolved conjectures are recorded
as such, rather than marked ``refuted'' from the absence of a proof.

\section{D\'erivateurs, localization, and information forgotten}
\label{sec:root-derivators}

The supplied DER\_LOCALIZERS intake identifies the linear
``Thm(?)'' at cote 134-8/048--054 and its remaining acyclicity
condition C.10.  Its original scans and the full proof of its
general counterexample C.26 are not supplied.  A specific motivic
instance is nevertheless established in the literature:
Voevodsky (2000), Theorem 3.1.12, gives strict homotopy invariance
for Nisnevich cohomology of homotopy-invariant presheaves with
transfers over a perfect field; Proposition 3.2.3 identifies the
Suslin-complex localization; Theorem 3.2.6 gives the embedding.
Those are geometric inputs to the localization theorem, independent
of RH.

\paragraph{Finding BR-DER-01 (constant rational coefficients detect exactly
rational homology).}
Let $\mathcal C$ be a nonzero presentable stable
$\mathbb Q$-linear category, with a specified nonzero object
$\mathbf1$ and the colimit-preserving exact action
$D(\mathbb Q)\to\mathcal C$, $V\mapsto V\otimes\mathbf1$.
For a functor $f:A\to B$ between small categories, the induced map
on homotopy colimits of the constant diagram $\mathbf1$ is an
equivalence if and only if $Nf$ is a rational homology equivalence.
Thus this particular constant-object localizer does not change
when one equips $\mathcal C$ with additional Tate twists,
Frobenius actions, or $L$-functions while retaining the specified
rational action and a nonzero unit.

\begin{proof}
The tensoring of a presentable stable category over spaces, followed
by the $\mathbb Q$-linear action, identifies the homotopy colimit
of a constant diagram on $A$ with
$C_*(NA;\mathbb Q)\otimes\mathbf1$.  Let $K$ be the cone in
$D(\mathbb Q)$ of the chain map induced by $Nf$.  Exactness shows
that the homotopy-colimit map is an equivalence exactly when
$K\otimes\mathbf1=0$.
Every complex of vector spaces splits as its homology together
with a contractible complex: choose complements to boundaries
inside cycles and to cycles inside each chain group.  The
differential identifies the latter complement with the
boundaries one degree lower, giving the contractible summands.
This works for unbounded complexes degree by degree.  Hence
$K\simeq\bigoplus_i H_i(K)[i]$ in $D(\mathbb Q)$.
If some $H_i(K)\ne0$, choose a nonzero vector and a linear
projection onto its span.  This exhibits $\mathbb Q[i]$ as a
retract of $K$.  Its image $\mathbf1[i]$ is then a nonzero
retract of $K\otimes\mathbf1$, a contradiction.  Conversely,
$K=0$ plainly implies $K\otimes\mathbf1=0$.
Thus the latter vanishing is equivalent to zero rational homology
of the cone, which is exactly the asserted condition on $Nf$.
\end{proof}

The quantifiers are deliberately explicit.  This proves the
constant-unit statement needed for rational motivic categories;
it does not identify all coefficient-dependent localizers, and
does not assert that every reversible localizer is of this form.
It also avoids the intake's unverified blanket full-faithfulness
assertion for arbitrary coefficient rings.

\paragraph{Finding BR-DER-02 (the reduced Tate object is essential).}
In geometric motives over a field, with the rational or usual
integral coefficients for which the cited motivic identities hold,
\[
 M(\mathbb G_m)\simeq\mathbf1\oplus\mathbf1(1)[1],\qquad
 \widetilde M(\mathbb G_m)[-1]\simeq\mathbf1(1).
\]
The formula in the DER\_LOCALIZERS intake omits the reduction if
$\mathbb Z_{\mathrm{tr}}(\mathbb G_m)$ means the full representable
presheaf with transfers.  Its full localized motive shifted by
$-1$ has an additional $\mathbf1[-1]$ summand.

\begin{proof}
Cover $\mathbb P^1$ by its two affine lines.  The motivic
Mayer--Vietoris triangle is
\[
 M(\mathbb G_m)\longrightarrow\mathbf1\oplus\mathbf1
 \longrightarrow M(\mathbb P^1)\longrightarrow M(\mathbb G_m)[1].
\]
The projective-line decomposition is
$M(\mathbb P^1)=\mathbf1\oplus\mathbf1(1)[2]$.
Each affine-line inclusion has degree-one component to
$\mathbf1$ and zero component to the Tate summand; the latter
lies in $\operatorname{Hom}(\mathbf1,\mathbf1(1)[2])
=\operatorname{Pic}(\operatorname{Spec}k)=0$.
After an invertible change of basis in $\mathbf1\oplus\mathbf1$,
the middle map is the identity on one unit summand and zero on
the other, with zero map to $\mathbf1(1)[2]$.  Taking its fibre
gives $M(\mathbb G_m)=\mathbf1\oplus\mathbf1(1)[1]$.
The rational point $1\in\mathbb G_m$ splits the augmentation;
its reduced fibre is exactly the second summand.  Shifting proves
the assertion.  The projective-line definition and reduction are
Voevodsky (2000), \S2.1 following Proposition 2.1.3; the Picard
calculation is Corollary 3.4.3.  This proof spells out the corrected
form rather than reproving the construction of motives.
\end{proof}

\section{A concrete boundary for the anabelian and Teichm\"uller bridges}
\label{sec:root-anabelian}

\paragraph{Finding BR-ANA-01 (identical abelian cohomological data do not
reconstruct a four-punctured sphere).}
The two hyperbolic curves over $\mathbb Q$
\[
 U_2=\mathbb P^1\setminus\{0,1,\infty,2\},\qquad
 U_3=\mathbb P^1\setminus\{0,1,\infty,3\}
\]
are not isomorphic even over $\overline{\mathbb Q}$, although
their geometric $\ell$-adic cohomology representations are
isomorphic in every degree.  For every prime $p\notin\{2,3\}$
and integer $r\geq1$ they have the same point count and local zeta
function:
\[
 \#U_i(\mathbb F_{p^r})=p^r-3,\qquad
 Z(U_i/\mathbb F_p,T)=\frac{(1-T)^3}{1-pT}.
\]
Thus the global good-prime cohomological zeta functions agree,
even after imposing RH on their Dirichlet/Tate factors.

\begin{proof}
An isomorphism of the affine smooth curves extends uniquely to
their smooth projective compactifications and permutes their
four missing points.  The six possible cross-ratios of an
unordered quadruple with parameter $t$ are
\[
 t,\quad 1-t,\quad t^{-1},\quad(1-t)^{-1},\quad
 t/(t-1),\quad(t-1)/t.
\]
For $t=2$ the resulting set is $\{2,-1,1/2\}$; for $t=3$ it
is $\{3,-2,1/3,-1/2,3/2,2/3\}$.  They are disjoint, so no
projective-linear transformation carries one unordered boundary
to the other.  This proves nonisomorphism.

For either boundary $D$ the localization sequence on
$\mathbb P^1_{\overline{\mathbb Q}}$ is Galois equivariant.
Since its four points are rational, the residue sequence gives
\[
 H^0(U_i,\mathbb Q_\ell)=\mathbb Q_\ell,\qquad
 H^1(U_i,\mathbb Q_\ell)
 =\ker\bigl(\mathbb Q_\ell(-1)^4
             \xrightarrow{\text{sum}}\mathbb Q_\ell(-1)\bigr)
 \simeq\mathbb Q_\ell(-1)^3,
\]
and higher ordinary cohomology vanishes.  This determines the
representations independently of $t$.  Alternatively compactly
supported cohomology has $H_c^1=\mathbb Q_\ell^3$ and
$H_c^2=\mathbb Q_\ell(-1)$.
At each stated prime the four rational points remain distinct,
so subtracting them from $\mathbb P^1(\mathbb F_{p^r})$ gives
$p^r+1-4$.  Exponentiating
$\sum_{r\geq1}(p^r-3)T^r/r$ yields the claimed zeta function.
\end{proof}

This is an explicit counterexample to reconstructing the
Teichm\"uller parameter from the indicated cohomological or
$L$-function data.  It does not contradict Grothendieck's
anabelian reconstruction proposals, whose input is a nonabelian
arithmetic fundamental group with its Galois action and peripheral
structure.  It shows why the abelian Artin--Tate result, even with
effective recovery of all its finite characters, is not itself
that nonabelian input.  A comparison with $\widehat{\mathrm{GT}}$
would have to specify and prove compatibility of a construction
with the noncommutative fundamental groups, their inertia, and
the tower maps.  None is constructed by the coincident words
``hexagon'' and ``tower'' in the supplied proof sketch.

